\phantomsection\bheading{section}{10.\quad Candidate language changes, for the owner}{section}{\protect\numberline{10.}Candidate language changes, for the owner}

The owner has said the specification is open where the checker needs it. Each change below says why, what precision it buys, what it costs developers, and how it weighs against rule 7 (“guarantees, not restrictions”). None is required for soundness; each is a trade.

\textbf{Change 1 — widen the demand set to the same-shape rebuilds.} Make \code{map}, \code{indexedMap}, \code{filter}, \code{filterMap}, \code{reverse}, \code{sort}, \code{sortBy}, \code{sortWith}, \code{take} and \code{slice} consuming writers: \code{map} is \code{a[i] = f(a[i])}, \code{filter} compacts in place, \code{sort} is \code{a.sort(cmp)} (stable in every current engine), \code{take n} is \code{a.length = n}. \emph{Reason:} the owner's programs edit lists this way, not with \code{set}/\code{push} (§8): without it the rule changes nothing in TodoMVC or Conduit, and nested-list edits through \code{map} are rejected (§5.4). \emph{Buys:} in-place \code{Toggle}, \code{Destroy}, \code{ClearCompleted}; unique elements inside \code{map}'s and \code{update}'s callbacks where the contents are exclusive (§2.9), and the trie's last reason to exist gone. \emph{Costs:} \code{List.map xs f} where \code{xs} is read later is now an error; today's identity promise “\code{map} returns \code{xs} itself when every result is \code{===}” must go for \code{map} (it would be true by construction — the array is \code{xs} — but the renderer's A4 obligation widens: a \code{For}'s \code{each} over a mapped list keeps its identity after an edit, so the identity walk must compare items, which \code{browser-direct} §6.2 already does). \emph{The wall:} the same \code{List.\allowbreak{}filter model.\allowbreak{}todos …} appears in \code{view}, where \code{browser-direct} compiles it into a derived value recomputed while the page holds \code{model}, and in \code{verifyAll}'s re-evaluation (§5.9); there it cannot be a demand. So change 1 is either \textbf{two functions} (\code{List.map} fresh, \code{List.edit}/\code{List.keep} in place — FP²'s duplication, §1.4.1) or \textbf{one spelling decided by position} (a demand in an \code{update} arm, a fresh build in a \code{view} or a derived value), which is a static rule the owner can read but is not “one \code{map}”. Neither is a silent fallback; the second is the smaller surface. \emph{Rule 7:} it rejects more correct programs (every read-after-\code{map} in \code{update}) to buy a cost guarantee on the most common list operations. Recommended \textbf{only together with the report mode's numbers} for \code{rebuilds}, since it is where most false errors would come from, and only in the by-position form.

\textbf{Change 2 — turn the identity promises into summaries, or drop the ones that alias.} \code{backend.md} §4 promises \code{filter}, \code{map}, \code{slice}, \code{take}, \code{xs ++ []} return their input itself in some cases. Under the rule a result that \emph{may} be the input is $\mathit{res}\ni \pi _{1}$, so the input must be unique at any later write of the result. \emph{Reason:} otherwise \code{xs ▷ List.\allowbreak{}filter p ▷ List.\allowbreak{}push y} demands \code{xs} unique though \code{filter} usually returns a fresh array. \emph{Buys:} fresh results everywhere, so pipelines never demand their source. The renderer's skip on \code{===} for a \code{filter} that kept everything is lost (one O(n) compare instead of one pointer test). \emph{Costs:} a measurable per-element compare in the renderer for lists whose \code{filter} keeps all; a golden change in \code{run/ListIdentity}. \emph{Rule 7:} no restriction either way; a representation choice. Recommended: drop the promises for \code{filter}, \code{map}, \code{slice}, \code{take}, \textbf{\code{[] ++ ys}} (today \code{append} returns \code{ys} itself when \code{xs} is empty, so a consumed \code{append} would hand out \code{ys}'s cell as its result) and \textbf{\code{concat}'s “the one non-empty list itself”} — they become always-fresh, or in-place under change 1 — and keep them for out-of-range \code{set}/\code{swap}/\code{removeAt} (no-ops on the unique array). \code{backend.md} §4 and §8 already disagree on \code{filter} (“returns \code{xs} itself when every element is kept” against “a new list and still is, \code{run/ListIdentity}”); the specification must settle it before $\mathit{res}$ can be declared.

\textbf{Change 3 — \code{List.\allowbreak{}copy : List a → List a}.} New, shallow, $\mathsf{fresh}$. \emph{Reason:} the explicit escape hatch; the machine-applicable fix. \emph{Costs:} none; a \code{copy\_\allowbreak{}of\_\allowbreak{}unique} \textbf{warning} (never an error) for a copy the checker can see is not needed. \emph{Rule 7:} an escape hatch, which rule 7 asks for.

\textbf{Change 4 — \code{++} on lists consumes its left operand.} \code{xs ++ ys} lowers to \code{List.append} (O6), which under the rule is \code{push} each of \code{ys} onto \code{xs}: \code{xs} consumed, \code{ys} borrowed. \emph{Reason:} consistency with \code{[ …xs, …ys ]}. \emph{Costs:} \code{xs ++ ys} with \code{xs} read later is an error; \code{Basics.append} over \code{appendable} stays a copy. TodoMVC's \code{model.\allowbreak{}todos ++ [ new ]} becomes a push. \emph{Rule 7:} same shape as change 1, smaller. Recommended with change 1.

\textbf{Change 5 — ownership words on \code{foreign} declarations, and optionally on \code{pub} annotations.} A \code{foreign} parameter may say $\mathsf{consumed}$, $\mathsf{borrowed}$ (the default) or $\mathsf{shared}$; a result $\mathsf{fresh}$ (the default) or $=\pi _{i}$. A \code{pub} beni declaration \emph{may} say the same in its annotation, and when it does the inferred summary must be at most what is written (an annotation can only claim \emph{less} ownership than the body takes — $\mathsf{borrowed}$ where the body consumes is an error at the declaration; $\mathsf{consumed}$ where the body borrows is accepted and freezes the interface at $\mathsf{consumed}$). \emph{Reason:} \code{foreign} has no body; platform contracts need the word (A2, A6); and \code{core/List}'s writers are beni over \code{Js}, whose bodies cannot show what they promise (§2.6), so \textbf{core and the platforms may assert a summary the body does not prove} — trusted, like the rung — while every other package may only \emph{narrow} an inferred one. For \code{pub} code it is Klabnik's and OxCaml's answer to interface churn: an annotation makes the summary a declared contract, so a body edit cannot change it (research 19 §4: annotated declarations measured zero interface changes). The vocabulary must also reach a function type inside a platform record — \code{Tea.sandbox}'s \code{\{ update : sync (msg, model → model), … \}} field is where P1's “\code{model} supplied owned” is written, as \code{sync} already is. \emph{Buys:} stable interfaces for library authors who want them; the pessimistic default for platform \code{foreign}s without words. \emph{Costs:} a vocabulary users may write but never must — the owner asked for no annotations in user code, and this keeps that: inference is complete without them. \emph{Rule 7:} an opt-in, not a restriction. Recommended for \code{foreign} (required for soundness anyway); optional for \code{pub}.

\textbf{Change 6 — closures that consume a capture are $\mathit{once}$, and a $\mathsf{many}$ position refuses them.} Already a consequence of the rules (§5.1); listed because it is a language-visible restriction: \code{List.\allowbreak{}map xs (λx → List.\allowbreak{}push acc x)} is an error whose fix is a fold. \emph{Buys:} soundness (two calls would write one array). \emph{Costs:} a shape some people write; the message gives the fold. \emph{Rule 7:} buys a correctness guarantee (no observable write), so it stays an error.

\textbf{Change 7 — a persistent vector as a library.} \code{core/Vector} (or \code{Persistent}), the E1tp trie written in beni with no in-place writes, for programs that keep versions. \emph{Reason:} the trie's users (undo, time travel, shared tails) lose their O(log n) versions when \code{List} stops being persistent. \emph{Costs:} a second sequence type in the library, which the owner removed from the \emph{language} (“one sequence type”, W35); as a library module it is \code{Dict}'s status, not \code{List}'s. \emph{Rule 7:} a capability filled inside the wall rather than a “work around it”.

\textbf{Change 8 — a deque headroom for prepend, or none.} \code{[ x, …xs ]} on a unique plain array is \code{unshift}, O(n). A unique view with headroom ($o>0$) could take an O(1) prepend by writing \code{b[o − 1]}; \code{cons} on a plain array could allocate headroom on its first conversion. \emph{Reason:} Elm-style prepend accumulators (research 46 §11's rows). \emph{Buys:} linear prepend loops without the building-loop rule. \emph{Costs:} a fourth list form or a growth policy in \code{core/List.js}; measured, not assumed. \emph{Rule 7:} a cost, not a guarantee; decide on §9.3's \code{prepends-loose} count.

\textbf{Change 9 — \code{Send msg} consumes its message.} A fiber that sends a value and reads it after is an error at the sender (P6). \emph{Buys:} unique payloads at the arm. \emph{Costs:} \code{send (Got rows); use rows} becomes \code{send (Got (List.\allowbreak{}copy rows)); use rows}. \emph{Rule 7:} the alternative is every payload shared at every arm; this is the cheaper default.

\textbf{Change 10 — evaluation-order dependence, stated once more.} The rule depends on \code{language.md} §6's order being the order the emitted code runs in; it already is, and A5 says the optimiser keeps it. No change, but the dependence should be written into §6 (“in-place writes are calls”).
